\scanpage{10}\mspage{7}

On a aussi l'hom $\Lambda$ de degré $-2$, défini par
\[
 \left\{\ \cdots\cdots\ \right. \tag*{(6.5)}
\]

\uline{Théorème} Les conditions suivantes \inserted{(i) à (iii)} sont équivalentes

(i) Les conditions \inserted{$(6.1)_i$ et} $(6.2)_i$ pour tt $i$, et les $\pi_i\in \unsure{E|\mathcal E}$.

\deleted{[\unsure{La graduation} de $E^*$ peut être définie par des $\pi_i$, et pour un \deleted{\unsure{quel}} \unsure{tel choix}]} \deleted{(\unsure{ii}) \unsure{la condition} $(6.3)_i$ pour tt $i$ [avec $\pi_i$ remplacés par $\pi_i$ \unsure{\ldots}]} \deleted{(\unsure{ii})}

\inserted{(ii) \deleted{\unsure{les} \illegible{1} et} On a la condition $(6.3)_i$ pour tout $i$}

(iii) \deleted{\unsure{L'hom}} \inserted{On a} \inserted{$(6.1)_i$ pour tt $i$, (d'où $\Lambda$) et}
\[
 \underline{\Lambda\in E^{(-2)}}\tag*{(6.6)}
\]

\unsure{(iv)} Ces conditions impliquent de plus

(iv) Les $\pi_{i,\alpha}$ sont dans $E$, \deleted{les $\pi_i$ sont dans $E$.}

(v) Les $\pi_i\in E$, \inserted{et ils} sont définis \inserted{(en termes intrinsèques à $E$ gradué)} comme formant \deleted{les seuls projecteurs} la seule famille d'élémts de $E$ satisfaisant (1.1.) et (1.3) (a) ou b), au choix \ldots, \deleted{\unsure{E}} ou encore (1.4)\ldots).

\deleted{(vi) L'élèmt $\Lambda$ de $E^{(-2)}$ est uniquement déterminé, \deleted{\unsure{par les}} (en termes de \inserted{l'anneau gradué} $E^*$ et l'élèmt $L\in E^{(\unsure{2|1})}$, par les conditions suivantes :
\[
 \left\{\begin{aligned}
 L\Lambda&=1-\sum_{0\le i\le n}\pi_{i,0}\\
 \Lambda L&=1-\sum_{n\le i\le 2n}\pi_{i,0}
 \end{aligned}\right.
 \tag*{(6.7)}
\]}

\uline{Dém} L'équivalence de (i) et (ii) est claire, et le fait que c'est impliqué par (iii) (en prenant $w_i=\Lambda^{\deleted{\unsure{(}}n-i\deleted{\unsure{)}}}$). \deleted{\unsure{Il}} \inserted{Reste à} prouver que (i), (ii) implique (iii) à (\unsure{vi|v}). On \deleted{\unsure{sait déjà}} \inserted{a déjà signalé} que ça implique \deleted{\unsure{(6.3)}} (v).
% archive sheet number: 9
